\documentclass{article}

% --- Paquetes para Estilo y Matemáticas ---
\usepackage[utf8]{inputenc}
\usepackage[T1]{fontenc}
\usepackage{amsmath} % Para notación matemática y entornos de ecuación
\usepackage{array} % Para control de tablas
\usepackage{booktabs} % Para tablas profesionales (sin líneas verticales)
\usepackage{caption} % Para captions de tablas
\usepackage{geometry} % Para márgenes
\usepackage{tcolorbox} % Para recuadros
\usepackage{graphicx} % Para incluir imágenes
\usepackage{fancyhdr} % Para encabezados y pies de página
\usepackage{lastpage} % Para número de páginas totales
\usepackage{titlesec} % Para personalizar secciones
\usepackage{enumitem} % Para personalizar listas
\usepackage{multicol} % Para columnas múltiples

% --- Configuración de Estilo ---
\geometry{a4paper, margin=1in, headheight=15pt}

% --- Configuración de Encabezados y Pies ---
\pagestyle{fancy}
\fancyhf{}
\fancyhead[L]{\small\textbf{Compendio de Cálculo}}
\fancyhead[C]{\small Derivadas e Integrales}
\fancyhead[R]{\small\thepage/\pageref{LastPage}}
\fancyfoot[C]{\footnotesize Resumen para estudiantes de matemáticas}
\renewcommand{\headrulewidth}{0.4pt}
\renewcommand{\footrulewidth}{0.2pt}

% --- Configuración de Títulos ---
\titleformat{\section}
{\normalfont\Large\bfseries}
{\thesection}{1em}{}

\titleformat{\subsection}
{\normalfont\large\bfseries}
{\thesubsection}{0.5em}{}

% --- Redefinición de Listas ---
\setlist[itemize]{leftmargin=*, itemsep=2pt, parsep=2pt}
\setlist[enumerate]{leftmargin=*, itemsep=2pt, parsep=2pt}

% --- Comandos Personalizados ---
\newcommand{\highlight}[1]{\textbf{#1}}
\newcommand{\formula}[1]{\ensuremath{#1}}
\newcommand{\deriv}[2]{\frac{d#1}{d#2}}
\newcommand{\partialderiv}[2]{\frac{\partial #1}{\partial #2}}
\newcommand{\integral}[2]{\int #1 \, d#2}
\newcommand{\examplebox}[1]{\begin{tcolorbox}[colback=gray!5!white, colframe=black, fonttitle=\bfseries, title=Ejemplo, #1\end{tcolorbox}}

% --- Portada ---
\title{\Huge\textbf{Compendio Completo de Cálculo}\\
\Large Derivadas e Integrales: Fórmulas, Propiedades y Técnicas}
\author{\Large Departamento de Matemáticas\\[0.5em]
\small Última actualización: \today}
\date{}

\begin{document}

% Portada
\begin{titlepage}
    \maketitle
    \centering
    \vspace{2cm}
    
    \begin{tcolorbox}[
        enhanced,
        sharp corners,
        colback=white,
        colframe=black,
        width=0.8\textwidth,
        center,
        fontupper=\large,
        drop shadow={black!50!white}
    ]
        \centering
        \textbf{Contenido del Documento}\\[0.5em]
        \begin{minipage}{0.9\textwidth}
            \begin{itemize}
                \item Sección 1: Propiedades fundamentales y teoremas
                \item Sección 2: Derivadas - Fórmulas y técnicas
                \item Sección 3: Integrales - Fórmulas y técnicas
                \item Sección 4: Clasificación y conceptos avanzados
                \item Sección 5: Glosario y referencias
                \item Apéndice A: Ejemplos resueltos
                \item Apéndice B: Tabla paralela de 40 derivadas e integrales
            \end{itemize}
        \end{minipage}
    \end{tcolorbox}
    
    \vfill
    \rule{0.8\textwidth}{0.4pt}\\[0.5em]
    \small Este documento resume las reglas fundamentales del cálculo diferencial e integral,\\
    incluyendo técnicas necesarias para resolver diferentes tipos de problemas.
\end{titlepage}

% Índice
\tableofcontents
\newpage

% ============================================
% SECCIÓN 1: FUNDAMENTOS
% ============================================
\section{Fundamentos del Cálculo}

\subsection{Propiedades Lineales y Operaciones Básicas}

\begin{multicols}{2}
\noindent
\textbf{Derivadas:}
\begin{itemize}
    \item \formula{\deriv{}{x}[c] = 0}
    \item \formula{\deriv{}{x}[c \cdot f(x)] = c \cdot f'(x)}
    \item \formula{\deriv{}{x}[f(x) \pm g(x)] = f'(x) \pm g'(x)}
    \item \formula{\deriv{}{x}[f(x) \cdot g(x)] = f'(x)g(x) + f(x)g'(x)}
    \item \formula{\deriv{}{x}\left[\frac{f(x)}{g(x)}\right] = \frac{f'(x)g(x) - f(x)g'(x)}{[g(x)]^2}}
    \item \formula{\deriv{}{x}[f(g(x))] = f'(g(x)) \cdot g'(x)}
\end{itemize}

\columnbreak

\noindent
\textbf{Integrales:}
\begin{itemize}
    \item \formula{\integral{c \cdot f(x) \pm g(x)}{x} = c \cdot \integral{f(x)}{x} \pm \integral{g(x)}{x}}
    \item \formula{\integral{0}{x} = C}
    \item \formula{\int_a^b f(x) dx = - \int_b^a f(x) dx}
    \item \formula{\int_a^c f(x) dx = \int_a^b f(x) dx + \int_b^c f(x) dx}
    \item \formula{\int_a^b f(x) dx = F(b) - F(a)}
\end{itemize}
\end{multicols}

\subsection{Teorema Fundamental del Cálculo y Regla de Leibniz}

\begin{center}
\begin{tabular}{@{}lll@{}}
\toprule
\textbf{Caso} & \textbf{Integral} & \textbf{Derivada ($\frac{d}{dx}$)} \\
\midrule
Límite Superior Variable & \formula{\int_a^x f(t) dt} & \formula{f(x)} \\
Límite Variable Compuesto & \formula{\int_a^{g(x)} f(t) dt} & \formula{f(g(x)) \cdot g'(x)} \\
Ambos Límites Variables & \formula{\int_{h(x)}^{g(x)} f(t) dt} & \formula{f(g(x)) \cdot g'(x) - f(h(x)) \cdot h'(x)} \\
\bottomrule
\end{tabular}
\end{center}

% ============================================
% SECCIÓN 2: DERIVADAS
% ============================================
\section{Derivadas: Fórmulas y Técnicas}

\subsection{Funciones Elementales}

\begin{center}
\small
\begin{tabular}{@{}lll@{}}
\toprule
\textbf{Función} & \textbf{Derivada} & \textbf{Ejemplo} \\
\midrule
\formula{x^n} & \formula{n x^{n-1}} & \formula{(x^5)' = 5x^4} \\
\formula{u^n} & \formula{n u^{n-1} \cdot u'} & \formula{((3x^2)^5)' = 5(3x^2)^4 \cdot 6x} \\
\formula{e^u} & \formula{e^u \cdot u'} & \formula{(e^{4x})' = 4e^{4x}} \\
\formula{a^u} & \formula{a^u \ln(a) \cdot u'} & \formula{(2^{\sin x})' = 2^{\sin x}\ln(2)\cos x} \\
\formula{\ln u} & \formula{\frac{u'}{u}} & \formula{(\ln(x^3+1))' = \frac{3x^2}{x^3+1}} \\
\formula{\log_a u} & \formula{\frac{u'}{u \ln a}} & \formula{(\log_{10}(2x))' = \frac{1}{x\ln 10}} \\
\formula{\sin u} & \formula{\cos u \cdot u'} & \formula{(\sin(5x))' = 5\cos(5x)} \\
\formula{\cos u} & \formula{-\sin u \cdot u'} & \formula{(\cos(x^2))' = -2x\sin(x^2)} \\
\formula{\tan u} & \formula{\sec^2 u \cdot u'} & \formula{(\tan(\sqrt{x}))' = \frac{\sec^2(\sqrt{x})}{2\sqrt{x}}} \\
\bottomrule
\end{tabular}
\captionof{table}{Derivadas de funciones elementales}
\end{center}

\subsection{Funciones Trigonométricas Inversas}

\begin{center}
\small
\begin{tabular}{@{}ll@{}}
\toprule
\textbf{Función} & \textbf{Derivada} \\
\midrule
\formula{\arcsin u} & \formula{\frac{u'}{\sqrt{1 - u^2}}} \\
\formula{\arccos u} & \formula{-\frac{u'}{\sqrt{1 - u^2}}} \\
\formula{\arctan u} & \formula{\frac{u'}{1 + u^2}} \\
\formula{\text{arccot } u} & \formula{-\frac{u'}{1 + u^2}} \\
\formula{\text{arcsec } u} & \formula{\frac{u'}{|u|\sqrt{u^2 - 1}}} \\
\formula{\text{arccsc } u} & \formula{-\frac{u'}{|u|\sqrt{u^2 - 1}}} \\
\bottomrule
\end{tabular}
\captionof{table}{Derivadas de funciones trigonométricas inversas}
\end{center}

% ============================================
% SECCIÓN 3: INTEGRALES
% ============================================
\section{Integrales: Fórmulas y Técnicas}

\subsection{Integrales Inmediatas}

\begin{center}
\small
\begin{tabular}{@{}lll@{}}
\toprule
\textbf{Función} & \textbf{Integral} & \textbf{Ejemplo} \\
\midrule
\formula{x^n \ (n \neq -1)} & \formula{\frac{x^{n+1}}{n+1} + C} & \formula{\int x^3 dx = \frac{x^4}{4} + C} \\
\formula{\frac{1}{x}} & \formula{\ln|x| + C} & \formula{\int \frac{1}{2x} dx = \frac{1}{2}\ln|x| + C} \\
\formula{e^{ax}} & \formula{\frac{1}{a}e^{ax} + C} & \formula{\int e^{5x} dx = \frac{1}{5}e^{5x} + C} \\
\formula{\sin x} & \formula{-\cos x + C} & \formula{\int \sin(3x) dx = -\frac{1}{3}\cos(3x) + C} \\
\formula{\cos x} & \formula{\sin x + C} & \formula{\int \cos(2x) dx = \frac{1}{2}\sin(2x) + C} \\
\formula{\frac{1}{a^2 + x^2}} & \formula{\frac{1}{a}\arctan\left(\frac{x}{a}\right) + C} & \formula{\int \frac{1}{4+x^2} dx = \frac{1}{2}\arctan\left(\frac{x}{2}\right) + C} \\
\bottomrule
\end{tabular}
\captionof{table}{Fórmulas de integración básica}
\end{center}

\subsection{Técnicas de Integración}

\begin{center}
\small
\begin{tabular}{@{}lp{8cm}@{}}
\toprule
\textbf{Método} & \textbf{Descripción y Fórmula} \\
\midrule
\highlight{Sustitución} & Análogo a la Regla de la Cadena: \formula{\int f(g(x)) g'(x) dx = \int f(u) du} \\
\highlight{Integración por Partes} & Análogo a la Regla del Producto: \formula{\int u dv = uv - \int v du} \\
\highlight{Sustitución Trigonométrica} & Para formas con raíces: \\
& \quad \formula{\sqrt{a^2 - x^2} \rightarrow x = a\sin\theta} \\
& \quad \formula{\sqrt{a^2 + x^2} \rightarrow x = a\tan\theta} \\
& \quad \formula{\sqrt{x^2 - a^2} \rightarrow x = a\sec\theta} \\
\highlight{Fracciones Parciales} & Para funciones racionales \formula{\frac{P(x)}{Q(x)}} descomponer en fracciones simples \\
\bottomrule
\end{tabular}
\captionof{table}{Técnicas de integración}
\end{center}

% ============================================
% SECCIÓN 4: CLASIFICACIÓN
% ============================================
\section{Clasificación y Conceptos Avanzados}

\subsection{Clasificación de Derivadas}

\begin{center}
\small
\begin{tabular}{@{}lllp{4cm}@{}}
\toprule
\textbf{Tipo} & \textbf{Variables} & \textbf{Notación} & \textbf{Descripción} \\
\midrule
Ordinaria & \formula{y = f(x)} & \formula{\frac{dy}{dx}} & Tasa de cambio respecto a una variable \\
Parcial & \formula{z = f(x,y)} & \formula{\partialderiv{z}{x}} & Cambio al variar una variable manteniendo otras constantes \\
Direccional & \formula{z = f(x,y)} & \formula{D_{\mathbf{u}} f} & Cambio en dirección específica \formula{\mathbf{u}} \\
Implícita & \formula{F(x,y)=0} & \formula{\frac{dy}{dx}} & Para funciones no despejadas explícitamente \\
\bottomrule
\end{tabular}
\captionof{table}{Tipos de derivadas según contexto}
\end{center}

\subsection{Clasificación de Integrales}

\begin{center}
\small
\begin{tabular}{@{}lllp{4cm}@{}}
\toprule
\textbf{Tipo} & \textbf{Notación} & \textbf{Resultado} & \textbf{Rol} \\
\midrule
Indefinida & \formula{\int f(x) dx} & \formula{F(x) + C} & Familia de antiderivadas \\
Definida & \formula{\int_a^b f(x) dx} & Número real & Área neta bajo la curva \\
Impropia & \formula{\int_a^\infty f(x) dx} & Límite & Intervalo infinito o discontinuidades \\
De línea & \formula{\int_C f(x,y) ds} & Número real & Integración sobre curvas \\
\bottomrule
\end{tabular}
\captionof{table}{Tipos de integrales}
\end{center}

\subsection{Derivadas de Orden Superior}

\begin{center}
\small
\begin{tabular}{@{}lll@{}}
\toprule
\textbf{Orden} & \textbf{Notación} & \textbf{Interpretación} \\
\midrule
Primera & \formula{f'(x)} & Pendiente, velocidad instantánea \\
Segunda & \formula{f''(x)} & Concavidad, aceleración \\
Tercera & \formula{f'''(x)} & Tasa de cambio de la aceleración \\
Enésima & \formula{f^{(n)}(x)} & Derivada de orden \formula{n} \\
\bottomrule
\end{tabular}
\captionof{table}{Derivadas de orden superior}
\end{center}

% ============================================
% SECCIÓN 5: GLOSARIO Y REFERENCIAS
% ============================================
\section{Glosario y Recursos}

\begin{multicols}{2}
\subsection*{Términos Clave}
\begin{itemize}
    \item \highlight{Antiderivada}: Función \formula{F(x)} cuya derivada es \formula{f(x)}
    \item \highlight{Área Neta}: Área considerando signo (positiva arriba del eje, negativa abajo)
    \item \highlight{Constante de Integración}: \formula{C}, representa familia de soluciones
    \item \highlight{Integrando}: Función a integrar
    \item \highlight{ILATE}: Regla nemotécnica para integración por partes
    \item \highlight{Función Racional}: Cociente de polinomios \formula{P(x)/Q(x)}
    \item \highlight{Discontinuidad}: Punto donde función no es continua
\end{itemize}

\columnbreak

\subsection*{Reglas Mnemotécnicas}
\begin{itemize}
    \item \highlight{ILATE}: Inversa, Logarítmica, Algebraica, Trigonométrica, Exponencial
    \item \highlight{LIATE}: Logarítmica, Inversa, Algebraica, Trigonométrica, Exponencial
    \item \highlight{PET}: Potencia, Exponencial, Trigonométrica (para integración por partes)
\end{itemize}
\end{multicols}

\begin{tcolorbox}[
    enhanced,
    sharp corners,
    colback=gray!5!white,
    colframe=black,
    title=\textbf{Consejos Prácticos},
    fonttitle=\bfseries,
    drop shadow={black!30!white}
]
\begin{itemize}
    \item \textbf{Siempre verifica} tu derivada/integral derivando/integrando hacia atrás
    \item \textbf{No olvides} la constante de integración \formula{C} en integrales indefinidas
    \item \textbf{Practica} la identificación del método adecuado para cada problema
    \item \textbf{Usa notación clara} y muestra todos los pasos intermedios
    \item \textbf{Revisa dominios} especialmente en funciones logarítmicas y raíces
\end{itemize}
\end{tcolorbox}

\subsection*{Referencias y Recursos}
\begin{itemize}
    \item Stewart, J. (2015). \textit{Cálculo: Trascendentes tempranas}. Cengage Learning.
    \item Larson, R., \& Edwards, B. (2017). \textit{Cálculo}. McGraw-Hill.
    \item Khan Academy: Cálculo diferencial e integral
    \item MIT OpenCourseWare: Cálculo de una variable
\end{itemize}

% ============================================
% APÉNDICE: EJEMPLOS PRÁCTICOS
% ============================================
\newpage
\appendix
\section{Ejemplos Resueltos}

\subsection{Ejemplo 1: Derivación Implícita}
\begin{tcolorbox}[colback=gray!5!white, colframe=black]
\textbf{Problema}: Encuentra \formula{\frac{dy}{dx}} para \formula{x^2 + y^2 = 25}

\textbf{Solución}:
\begin{align*}
\frac{d}{dx}(x^2) + \frac{d}{dx}(y^2) &= \frac{d}{dx}(25) \\
2x + 2y \cdot \frac{dy}{dx} &= 0 \\
2y \cdot \frac{dy}{dx} &= -2x \\
\frac{dy}{dx} &= -\frac{x}{y}
\end{align*}
\end{tcolorbox}

\subsection{Ejemplo 2: Integración por Partes}
\begin{tcolorbox}[colback=gray!5!white, colframe=black]
\textbf{Problema}: Calcula \formula{\int x e^x dx}

\textbf{Solución} (usando ILATE):
\begin{align*}
\text{Sea } u = x &\quad dv = e^x dx \\
du = dx &\quad v = e^x \\
\int x e^x dx &= x e^x - \int e^x dx \\
&= x e^x - e^x + C \\
&= e^x(x - 1) + C
\end{align*}
\end{tcolorbox}

\subsection{Ejemplo 3: Sustitución Trigonométrica}
\begin{tcolorbox}[colback=gray!5!white, colframe=black]
\textbf{Problema}: Calcula \formula{\int \frac{1}{\sqrt{9-x^2}} dx}

\textbf{Solución}:
\begin{align*}
\text{Sea } x &= 3\sin\theta, \quad dx = 3\cos\theta d\theta \\
\int \frac{1}{\sqrt{9-x^2}} dx &= \int \frac{1}{\sqrt{9-9\sin^2\theta}} \cdot 3\cos\theta d\theta \\
&= \int \frac{3\cos\theta}{3\cos\theta} d\theta = \int 1 d\theta \\
&= \theta + C = \arcsin\left(\frac{x}{3}\right) + C
\end{align*}
\end{tcolorbox}

% ============================================
% APÉNDICE: TABLA PARALELA
% ============================================
\newpage % Comienza la nueva sección de la tabla en una nueva página (Sección 1/2)
\section{Tabla Paralela de Derivadas e Integrales Inmediatas (40 Fórmulas)}
\label{sec:tabla_paralela_completa_parte1}

\subsection*{Instrucciones}
Esta tabla presenta derivadas inmediatas (columna izquierda) y sus correspondientes integrales inmediatas (columna derecha), organizadas en paralelo. Se utiliza una sintaxis compacta para maximizar la visibilidad de todas las fórmulas.

\begin{center}
\captionof{table}{Formulario Paralelo de Derivadas e Integrales (Parte 1)}
\footnotesize % Fuente pequeña para mayor compacidad
\setlength{\tabcolsep}{4pt} % Reducir espacio entre columnas
\begin{tabular}{|p{0.48\textwidth}|p{0.48\textwidth}|}
\toprule
\textbf{DERIVADAS INMEDIATAS} ($\deriv{}{x}[f(x)]$) & \textbf{INTEGRALES INMEDIATAS} ($\integral{f(x)}{x}$) \\
\midrule
\multicolumn{2}{|l|}{\textbf{1. Trigonométricas Básicas}} \\
\midrule
$\dfrac{d}{dx}[\sin x] = \cos x$ & $\displaystyle\int \cos x \, dx = \sin x + C$ \\
$\dfrac{d}{dx}[\cos x] = -\sin x$ & $\displaystyle\int \sin x \, dx = -\cos x + C$ \\
$\dfrac{d}{dx}[\tan x] = \sec^2 x$ & $\displaystyle\int \sec^2 x \, dx = \tan x + C$ \\
$\dfrac{d}{dx}[\cot x] = -\csc^2 x$ & $\displaystyle\int \csc^2 x \, dx = -\cot x + C$ \\
$\dfrac{d}{dx}[\sec x] = \sec x \tan x$ & $\displaystyle\int \sec x \tan x \, dx = \sec x + C$ \\
$\dfrac{d}{dx}[\csc x] = -\csc x \cot x$ & $\displaystyle\int \csc x \cot x \, dx = -\csc x + C$ \\
\midrule
\multicolumn{2}{|l|}{\textbf{2. Exponenciales}} \\
\midrule
$\dfrac{d}{dx}[e^x] = e^x$ & $\displaystyle\int e^x \, dx = e^x + C$ \\
$\dfrac{d}{dx}[e^{kx}] = ke^{kx}$ & $\displaystyle\int e^{kx} \, dx = \frac{1}{k}e^{kx} + C$ \\
$\dfrac{d}{dx}[a^x] = a^x \ln a$ & $\displaystyle\int a^x \, dx = \frac{a^x}{\ln a} + C$ \\
\midrule
\multicolumn{2}{|l|}{\textbf{3. Logarítmicas}} \\
\midrule
$\dfrac{d}{dx}[\ln|x|] = \frac{1}{x}$ & $\displaystyle\int \frac{1}{x} \, dx = \ln|x| + C$ \\
$\dfrac{d}{dx}[\ln|f(x)|] = \frac{f'(x)}{f(x)}$ & $\displaystyle\int \frac{f'(x)}{f(x)} \, dx = \ln|f(x)| + C$ \\
$\dfrac{d}{dx}[\log_a x] = \frac{1}{x \ln a}$ & $\displaystyle\int \frac{1}{x \ln a} \, dx = \log_a |x| + C$ \\
\midrule
\multicolumn{2}{|l|}{\textbf{4. Potenciales/Algebraicas}} \\
\midrule
$\dfrac{d}{dx}[x^n] = nx^{n-1}$ & $\displaystyle\int x^n \, dx = \frac{x^{n+1}}{n+1} + C, \quad n \neq -1$ \\
$\dfrac{d}{dx}[\sqrt{x}] = \frac{1}{2\sqrt{x}}$ & $\displaystyle\int \frac{1}{2\sqrt{x}} \, dx = \sqrt{x} + C$ \\
$\dfrac{d}{dx}\left[\frac{1}{x}\right] = -\frac{1}{x^2}$ & $\displaystyle\int -\frac{1}{x^2} \, dx = \frac{1}{x} + C$ \\
$\dfrac{d}{dx}\left[\frac{1}{x^n}\right] = -\frac{n}{x^{n+1}}$ & $\displaystyle\int -\frac{n}{x^{n+1}} \, dx = \frac{1}{x^n} + C$ \\
\midrule
\multicolumn{2}{|l|}{\textbf{5. Trigonométricas Inversas}} \\
\midrule
$\dfrac{d}{dx}[\arcsin x] = \frac{1}{\sqrt{1-x^2}}$ & $\displaystyle\int \frac{1}{\sqrt{1-x^2}} \, dx = \arcsin x + C$ \\
$\dfrac{d}{dx}[\arccos x] = -\frac{1}{\sqrt{1-x^2}}$ & $\displaystyle\int -\frac{1}{\sqrt{1-x^2}} \, dx = \arccos x + C$ \\
$\dfrac{d}{dx}[\arctan x] = \frac{1}{1+x^2}$ & $\displaystyle\int \frac{1}{1+x^2} \, dx = \arctan x + C$ \\
$\dfrac{d}{dx}[\operatorname{arccot} x] = -\frac{1}{1+x^2}$ & $\displaystyle\int -\frac{1}{1+x^2} \, dx = \operatorname{arccot} x + C$ \\
\midrule
\multicolumn{2}{|l|}{\textbf{6. Hiperbólicas}} \\
\midrule
$\dfrac{d}{dx}[\sinh x] = \cosh x$ & $\displaystyle\int \cosh x \, dx = \sinh x + C$ \\
$\dfrac{d}{dx}[\cosh x] = \sinh x$ & $\displaystyle\int \sinh x \, dx = \cosh x + C$ \\
$\dfrac{d}{dx}[\tanh x] = \operatorname{sech}^2 x$ & $\displaystyle\int \operatorname{sech}^2 x \, dx = \tanh x + C$ \\
$\dfrac{d}{dx}[\coth x] = -\operatorname{csch}^2 x$ & $\displaystyle\int -\operatorname{csch}^2 x \, dx = \coth x + C$ \\
\midrule
\multicolumn{2}{|l|}{\textbf{7. Funciones Compuestas Básicas}} \\
\midrule
$\dfrac{d}{dx}[\sin(kx)] = k\cos(kx)$ & $\displaystyle\int k\cos(kx) \, dx = \sin(kx) + C$ \\
$\dfrac{d}{dx}[\cos(kx)] = -k\sin(kx)$ & $\displaystyle\int -k\sin(kx) \, dx = \cos(kx) + C$ \\
$\dfrac{d}{dx}[\tan(kx)] = k\sec^2(kx)$ & $\displaystyle\int k\sec^2(kx) \, dx = \tan(kx) + C$ \\
$\dfrac{d}{dx}[e^{ax}] = ae^{ax}$ & $\displaystyle\int ae^{ax} \, dx = e^{ax} + C$ \\
$\dfrac{d}{dx}[\ln|ax|] = \frac{1}{x}$ & $\displaystyle\int \frac{1}{x} \, dx = \ln|ax| + C$ \\
\bottomrule
\end{tabular}
\end{center}

\newpage % <-- SALTO DE PÁGINA SOLICITADO

\section{Continuación: Radicales, Especiales y Funciones No Inmediatas}
\label{sec:tabla_paralela_continuacion}

\begin{center}
\captionof{table}{Formulario Paralelo de Derivadas e Integrales (Parte 2)}
\footnotesize % Fuente pequeña para mayor compacidad
\setlength{\tabcolsep}{4pt} % Reducir espacio entre columnas
\begin{tabular}{|p{0.48\textwidth}|p{0.48\textwidth}|}
\toprule
\textbf{DERIVADAS INMEDIATAS} ($\deriv{}{x}[f(x)]$) & \textbf{INTEGRALES INMEDIATAS} ($\integral{f(x)}{x}$) \\
\midrule
\multicolumn{2}{|l|}{\textbf{8. Con Radicales}} \\
\midrule
$\dfrac{d}{dx}[\sqrt{1-x^2}] = \frac{-x}{\sqrt{1-x^2}}$ & $\displaystyle\int \frac{-x}{\sqrt{1-x^2}} \, dx = \sqrt{1-x^2} + C$ \\
$\dfrac{d}{dx}[\sqrt{a^2-x^2}] = \frac{-x}{\sqrt{a^2-x^2}}$ & $\displaystyle\int \frac{-x}{\sqrt{a^2-x^2}} \, dx = \sqrt{a^2-x^2} + C$ \\
$\dfrac{d}{dx}[\sqrt{x^2+a^2}] = \frac{x}{\sqrt{x^2+a^2}}$ & $\displaystyle\int \frac{x}{\sqrt{x^2+a^2}} \, dx = \sqrt{x^2+a^2} + C$ \\
\midrule
\multicolumn{2}{|l|}{\textbf{9. Con Funciones Trigonométricas Especiales (Integrales de Funciones Trigonométricas)}} \\
\midrule
$\dfrac{d}{dx}[\ln|\sin x|] = \cot x$ & $\displaystyle\int \cot x \, dx = \ln|\sin x| + C$ \\
$\dfrac{d}{dx}[\ln|\cos x|] = -\tan x$ & $\displaystyle\int \tan x \, dx = -\ln|\cos x| + C$ \\
$\dfrac{d}{dx}[\ln|\sec x+\tan x|] = \sec x$ & $\displaystyle\int \sec x \, dx = \ln|\sec x+\tan x| + C$ \\
$\dfrac{d}{dx}[\ln|\csc x-\cot x|] = \csc x$ & $\displaystyle\int \csc x \, dx = \ln|\csc x-\cot x| + C$ \\
\midrule
\multicolumn{2}{|l|}{\textbf{10. Funciones que NO son Pares Perfectos (Ejemplos de Integrales por Partes)}} \\
\midrule
$\dfrac{d}{dx}[\arcsinh x] = \frac{1}{\sqrt{1+x^2}}$ & $\displaystyle\int \frac{1}{\sqrt{1+x^2}} \, dx = \operatorname{arsinh} x + C$ \\
$\dfrac{d}{dx}[\operatorname{arcosh} x] = \frac{1}{\sqrt{x^2-1}}$ & $\displaystyle\int \frac{1}{\sqrt{x^2-1}} \, dx = \operatorname{arcosh} x + C$ \\
$\dfrac{d}{dx}[x\ln x - x] = \ln x$ & $\displaystyle\int \ln x \, dx = x\ln x - x + C$ \\
$\dfrac{d}{dx}[x\sin x + \cos x] = x\cos x$ & $\displaystyle\int x\cos x \, dx = x\sin x + \cos x + C$ \\
$\dfrac{d}{dx}[x\cos x - \sin x] = -x\sin x$ & $\displaystyle\int x\sin x \, dx = -x\cos x + \sin x + C$ \\
$\dfrac{d}{dx}[xe^x - e^x] = xe^x$ & $\displaystyle\int xe^x \, dx = (x-1)e^x + C$ \\
$\dfrac{d}{dx}\left[\frac{x^2}{2}\arctan x - \frac{x}{2} + \frac{1}{2}\arctan x\right] = x\arctan x$ & $\displaystyle\int x\arctan x \, dx = \frac{x^2}{2}\arctan x - \frac{x}{2} + \frac{1}{2}\arctan x + C$ \\
$\dfrac{d}{dx}[x\arcsin x + \sqrt{1-x^2}] = \arcsin x$ & $\displaystyle\int \arcsin x \, dx = x\arcsin x + \sqrt{1-x^2} + C$ \\
$\dfrac{d}{dx}[x\arccos x - \sqrt{1-x^2}] = \arccos x$ & $\displaystyle\int \arccos x \, dx = x\arccos x - \sqrt{1-x^2} + C$ \\
\bottomrule
\end{tabular}
\end{center}

\subsection*{Explicación de la Organización}
\begin{itemize}
\item \textbf{Filas 1-30}: Pares perfectos donde la derivada de una función aparece directamente como integrando en la integral correspondiente.
\item \textbf{Filas 31-40}: Funciones donde la relación no es directa pero están relacionadas (integrales que se resuelven por partes o métodos especiales).
\item Cada par muestra la relación inversa entre derivación e integración.
\end{itemize}

\subsection*{Notas Importantes}
\begin{enumerate}
\item La constante de integración $C$ solo aparece en las integrales, no en las derivadas.
\item En algunos casos, las fórmulas de integrales tienen restricciones de dominio que no se muestran explícitamente.
\item Las funciones con valor absoluto ($|x|$, $|f(x)|$) son necesarias para cubrir todos los valores reales.
\item Las funciones hiperbólicas inversas tienen fórmulas logarítmicas equivalentes que no se muestran.
\end{enumerate}

\end{document}